\section{Method}\label{sec:3}

Given a sequence of human actions and a customizable world specification, our goal is to synthesize an egocentric video that reflects how a user navigates and interacts within the defined environment. 
To this end, AnchorWorld takes two types of control signals as input: embodied human motion for action control, and pose-associated anchor views for world customization.
We instantiate AnchorWorld with Wan~\cite{wan2025wan}, a flow-matching-based~\cite{lipman2022flow} DiT~\citep{peebles2023scalable} video generation model, and condition its video synthesis on the action and anchor-view signals.
The human motion is represented as a sequence of body actions derived from the SMPL-X parametric body model~\cite{pavlakos2019expressive}, denoted as
$M \in \sR^{f \times k \times 6}$, where $f$ is the number of frames and $k$ is the number of joints. 
Each joint state consists of its 3D position and 3D axis-angle rotation vector. 
The customizable world is defined by an initial egocentric view $I_0$ and a set of localized anchor views $\gS = \{ (\vI_i, \vc_i, \vt_i) \}_{i=1}^n$. 
Each anchor view contains an RGB image $\vI_i$, a 6-DoF viewpoint pose $\vc_i = [\mR_i \mid \vp_i] \in \sR^{3 \times 4}$, and an evolution prompt $\vt_i$ that describes the temporal change of local scene states.
Figure~\ref{fig:method} provides an overview of the proposed framework. 
We detail each component of our approach in the following subsections.



\section{Method}\label{sec:3}

Given a sequence of human actions and a customizable world specification, our goal is to synthesize an egocentric video that reflects how a user navigates and interacts within the defined environment. 
To this end, AnchorWorld takes two types of control signals as input: embodied human motion for action control, and pose-associated anchor views for world customization.
We instantiate AnchorWorld with Wan~\cite{wan2025wan}, a flow-matching-based~\cite{lipman2022flow} DiT~\citep{peebles2023scalable} video generation model, and condition its video synthesis on the action and anchor-view signals.
The human motion is represented as a sequence of body actions derived from the SMPL-X parametric body model~\cite{pavlakos2019expressive}, denoted as
$M \in \sR^{f \times k \times 6}$, where $f$ is the number of frames and $k$ is the number of joints. 
Each joint state consists of its 3D position and 3D axis-angle rotation vector. 
The customizable world is defined by an initial egocentric view $I_0$ and a set of localized anchor views $\gS = \{ (\vI_i, \vc_i, \vt_i) \}_{i=1}^n$. 
Each anchor view contains an RGB image $\vI_i$, a 6-DoF viewpoint pose $\vc_i = [\mR_i \mid \vp_i] \in \sR^{3 \times 4}$, and an evolution prompt $\vt_i$ that describes the temporal change of local scene states.
Figure~\ref{fig:method} provides an overview of the proposed framework. 
We detail each component of our approach in the following subsections.



\section{Method}\label{sec:3}

Given a sequence of human actions and a customizable world specification, our goal is to synthesize an egocentric video that reflects how a user navigates and interacts within the defined environment. 
To this end, AnchorWorld takes two types of control signals as input: embodied human motion for action control, and pose-associated anchor views for world customization.
We instantiate AnchorWorld with Wan~\cite{wan2025wan}, a flow-matching-based~\cite{lipman2022flow} DiT~\citep{peebles2023scalable} video generation model, and condition its video synthesis on the action and anchor-view signals.
The human motion is represented as a sequence of body actions derived from the SMPL-X parametric body model~\cite{pavlakos2019expressive}, denoted as
$M \in \sR^{f \times k \times 6}$, where $f$ is the number of frames and $k$ is the number of joints. 
Each joint state consists of its 3D position and 3D axis-angle rotation vector. 
The customizable world is defined by an initial egocentric view $I_0$ and a set of localized anchor views $\gS = \{ (\vI_i, \vc_i, \vt_i) \}_{i=1}^n$. 
Each anchor view contains an RGB image $\vI_i$, a 6-DoF viewpoint pose $\vc_i = [\mR_i \mid \vp_i] \in \sR^{3 \times 4}$, and an evolution prompt $\vt_i$ that describes the temporal change of local scene states.
Figure~\ref{fig:method} provides an overview of the proposed framework. 
We detail each component of our approach in the following subsections.



\section{Method}\label{sec:3}

Given a sequence of human actions and a customizable world specification, our goal is to synthesize an egocentric video that reflects how a user navigates and interacts within the defined environment. 
To this end, AnchorWorld takes two types of control signals as input: embodied human motion for action control, and pose-associated anchor views for world customization.
We instantiate AnchorWorld with Wan~\cite{wan2025wan}, a flow-matching-based~\cite{lipman2022flow} DiT~\citep{peebles2023scalable} video generation model, and condition its video synthesis on the action and anchor-view signals.
The human motion is represented as a sequence of body actions derived from the SMPL-X parametric body model~\cite{pavlakos2019expressive}, denoted as
$M \in \sR^{f \times k \times 6}$, where $f$ is the number of frames and $k$ is the number of joints. 
Each joint state consists of its 3D position and 3D axis-angle rotation vector. 
The customizable world is defined by an initial egocentric view $I_0$ and a set of localized anchor views $\gS = \{ (\vI_i, \vc_i, \vt_i) \}_{i=1}^n$. 
Each anchor view contains an RGB image $\vI_i$, a 6-DoF viewpoint pose $\vc_i = [\mR_i \mid \vp_i] \in \sR^{3 \times 4}$, and an evolution prompt $\vt_i$ that describes the temporal change of local scene states.
Figure~\ref{fig:method} provides an overview of the proposed framework. 
We detail each component of our approach in the following subsections.



\input{tex/figs/method}

\subsection{Hybrid-View Human Action Control}\label{sec:3.1}

\paragraph{Enhanced Egocentric Action Control via hybrid views.}
Human action contains rich spatial and interaction cues: the root trajectory determines global navigation, the limbs indicate potential interactions with the surrounding scene, and the head motion induces the egocentric viewpoint.
However, in first-person videos, most body parts are often outside the camera field of view, making direct supervision of full-body action control sparse and incomplete.
To overcome this limitation, we introduce third-person view (TPV) videos as auxiliary training data, where the full human body and its interactions with the surrounding scene are explicitly visible.
These videos provide rich interaction context and complete motion supervision, helping the model learn stronger spatial grounding between human motion and scene responses.
To support joint training on both TPV and first-person view (FPV) data within a unified framework, we formulate action conditioning in a projection-based manner.
Specifically, we represent the action condition by combining the full-body motion sequence with the camera trajectory, allowing the model to project 3D human motion into 2D visual observations under arbitrary viewpoints.
We first pre-train the model on large-scale and diverse TPV videos, where the camera parameters correspond to the external observation viewpoint, enabling the model to acquire projection knowledge and human-scene interaction priors.
Then, we adapt the model to egocentric simulation by aligning the camera parameters with the human head perspective in FPV data.
This design enables more accurate human-action control and stronger spatial pose awareness.

\paragraph{Spatial Pose Attention.} 
Inspired by prior work~\cite{fu20243dtrajmaster, li2025adaviewplanner}, we inject the pose conditions through a spatial pose attention mechanism. 
Specifically, a motion encoder first projects the input motion sequence $M \in \sR^{f \times k \times 6}$ into a latent embedding $\vz_m \in \sR^{f' \times k \times d}$, where $d$ is the model's hidden dimension. To ensure temporal alignment with the VAE-encoded~\cite{kingma2013auto} video latents, we employ temporal downsampling to match the temporal resolution $f'$. Analogously, a camera encoder processes the camera pose sequence $C \in \sR^{f \times 3 \times 4}$ into $\vz_c \in \sR^{f' \times 1 \times d}$, where the camera pose can represent either a third-person observation viewpoint or the first-person head viewpoint.

To exploit the inherent frame-wise correspondence between motion and video tokens, we concatenate the video tokens $\vz_v^{(t)}$ with the human motion tokens $\vz_m$ and camera pose tokens $\vz_c$ along the spatial dimension. This unified sequence is then processed by the spatial self-attention block:

\[
\begin{gathered}
\mT = [\vz_v^{(t)}; \vz_m; \vz_c] \in \sR^{f' \times (h \cdot w + k + 1) \times d}, \\
\vz_v^{(t)} = \vz_v^{(t)} + \text{Truncate}\left(\text{Attn}(\mW_Q \cdot \mT, \mW_K \cdot \mT, \mW_V \cdot \mT)\right)
\end{gathered}
\tag{1}
\]

The Truncate operator discards the auxiliary pose tokens, retaining only the updated video features.



\subsection{Evolvable Anchor-View Customization}\label{sec:3.2}

To enable evolvable world customization, we represent the environment with a set of anchor views. 
Each anchor view provides three types of localized world priors: an RGB image for visual appearance, a 3D pose for spatial grounding, and an evolution prompt for temporal state evolution. 

\paragraph{In-Context Anchor-View Priors.}
To incorporate anchor-view image priors while preserving the generative capability of the pre-trained video model, we adopt an in-context conditioning strategy~\cite{huang2026cinescene, ju2025fulldit, ye2025unic}. 
Specifically, the images of anchor views are encoded into latent tokens $\vz_s \in \sR^{f_s \times h \cdot w \times d}$, which are concatenated with the video latent tokens $\vz_v^{(t)} \in \sR^{f' \times h \cdot w \times d}$ along the frame dimension:
\[
    \mathcal{T}_{total} = [\vz_v^{(t)}; \vz_s] \in \sR^{(f' + f_s) \times h \cdot w \times d}. \tag{2}
\]

This design enables anchor views to guide world synthesis in-context, without requiring architectural modifications to the base model. We further employ 3D RoPE~\cite{su2024roformer} to differentiate anchor views by assigning them distinct frame-axis positions in the positional embedding space.

\input{tex/figs/train}


\paragraph{View Pose Injection.}
Since each view corresponds to a specific 3D location in the world, its spatial pose is essential for grounding the customized content. 
We therefore inject pose information for both generated video frames and anchor views. 
The camera poses are encoded into embeddings $\vz_{pose} \in \sR^{(f' + f_s) \times 1 \times d}$ and spatially broadcast to match the latent resolution, yielding $\vz_{pose} \in \sR^{(f' + f_s) \times h \cdot w \times d}$. 
Before the self-attention layers, the pose embeddings are added to the visual tokens:
\[
    \mathcal{T}_{total} = \mathcal{T}_{total} + \vz_{pose}. \tag{3}
\]

By coupling visual tokens with spatial poses, the model can distinguish anchor views located at different positions and associate the generated egocentric trajectory with the correct local constraints. 
% This pose-aware formulation helps maintain scene consistency when the user approaches, revisits, or moves across anchored regions.

\paragraph{Text-Driven Anchor-View Evolution.}
To enable dynamic world customization, each anchor view is paired with a localized evolution description $\vt_i$ that specifies its temporal scene changes.
We inject these descriptions through cross-attention, leveraging the semantic priors of the pre-trained video model.
To preserve the locality of dynamic instructions, we restrict the interaction between text prompts and visual tokens using an attention mask. 
For a text prompt $\vt_j$, its text keys are visible only to the generated video tokens and the corresponding anchor-view tokens $\vz_s^{(j)}$:
\[
    \mathcal{M}(q, k_j) = 
    \begin{cases} 
    0, & \text{if } q \in \vz_v \text{ or } q \in \vz_s^{(j)}, \\
    -\infty, & \text{otherwise}.
    \end{cases} \tag{4}
\]
This masked cross-attention enables anchor-specific text control, allowing local scene states to evolve over time while reducing interference across different anchor views.

\subsection{Progressive Multi-Stage Training Strategy}\label{sec:3.3}

To progressively equip the model with egocentric human action control and evolvable anchor-view customization, we adopt a multi-stage training strategy, as illustrated in Figure~\ref{fig:train}. 
%
\textbf{Stage I \& II: Hybrid-View Action Control Training.}
We train the model to learn action-conditioned generation from hybrid viewpoints, where TPV videos provide complete full-body motion supervision.
% and FPV videos adapt the model to egocentric simulation.
In Stage I, the model is trained on large-scale third-person videos, where the camera parameters represent external observation viewpoints.
In Stage II, we then adapt the model to first-person videos by aligning the camera trajectory with the head pose of the character.
%
\textbf{Stage III \& IV: Evolvable Anchor-View Customization Training.}
After establishing action controllability, we train the model to incorporate anchor-view priors for world customization.
In Stage III, we train the model on static scenes to learn pose-aware anchor-view conditioning for consistent egocentric roaming.
In Stage IV, we mix in dynamic data with evolution descriptions to model text-driven local state changes.






\subsection{Hybrid-View Human Action Control}\label{sec:3.1}

\paragraph{Enhanced Egocentric Action Control via hybrid views.}
Human action contains rich spatial and interaction cues: the root trajectory determines global navigation, the limbs indicate potential interactions with the surrounding scene, and the head motion induces the egocentric viewpoint.
However, in first-person videos, most body parts are often outside the camera field of view, making direct supervision of full-body action control sparse and incomplete.
To overcome this limitation, we introduce third-person view (TPV) videos as auxiliary training data, where the full human body and its interactions with the surrounding scene are explicitly visible.
These videos provide rich interaction context and complete motion supervision, helping the model learn stronger spatial grounding between human motion and scene responses.
To support joint training on both TPV and first-person view (FPV) data within a unified framework, we formulate action conditioning in a projection-based manner.
Specifically, we represent the action condition by combining the full-body motion sequence with the camera trajectory, allowing the model to project 3D human motion into 2D visual observations under arbitrary viewpoints.
We first pre-train the model on large-scale and diverse TPV videos, where the camera parameters correspond to the external observation viewpoint, enabling the model to acquire projection knowledge and human-scene interaction priors.
Then, we adapt the model to egocentric simulation by aligning the camera parameters with the human head perspective in FPV data.
This design enables more accurate human-action control and stronger spatial pose awareness.

\paragraph{Spatial Pose Attention.} 
Inspired by prior work~\cite{fu20243dtrajmaster, li2025adaviewplanner}, we inject the pose conditions through a spatial pose attention mechanism. 
Specifically, a motion encoder first projects the input motion sequence $M \in \sR^{f \times k \times 6}$ into a latent embedding $\vz_m \in \sR^{f' \times k \times d}$, where $d$ is the model's hidden dimension. To ensure temporal alignment with the VAE-encoded~\cite{kingma2013auto} video latents, we employ temporal downsampling to match the temporal resolution $f'$. Analogously, a camera encoder processes the camera pose sequence $C \in \sR^{f \times 3 \times 4}$ into $\vz_c \in \sR^{f' \times 1 \times d}$, where the camera pose can represent either a third-person observation viewpoint or the first-person head viewpoint.

To exploit the inherent frame-wise correspondence between motion and video tokens, we concatenate the video tokens $\vz_v^{(t)}$ with the human motion tokens $\vz_m$ and camera pose tokens $\vz_c$ along the spatial dimension. This unified sequence is then processed by the spatial self-attention block:

\[
\begin{gathered}
\mT = [\vz_v^{(t)}; \vz_m; \vz_c] \in \sR^{f' \times (h \cdot w + k + 1) \times d}, \\
\vz_v^{(t)} = \vz_v^{(t)} + \text{Truncate}\left(\text{Attn}(\mW_Q \cdot \mT, \mW_K \cdot \mT, \mW_V \cdot \mT)\right)
\end{gathered}
\tag{1}
\]

The Truncate operator discards the auxiliary pose tokens, retaining only the updated video features.



\subsection{Evolvable Anchor-View Customization}\label{sec:3.2}

To enable evolvable world customization, we represent the environment with a set of anchor views. 
Each anchor view provides three types of localized world priors: an RGB image for visual appearance, a 3D pose for spatial grounding, and an evolution prompt for temporal state evolution. 

\paragraph{In-Context Anchor-View Priors.}
To incorporate anchor-view image priors while preserving the generative capability of the pre-trained video model, we adopt an in-context conditioning strategy~\cite{huang2026cinescene, ju2025fulldit, ye2025unic}. 
Specifically, the images of anchor views are encoded into latent tokens $\vz_s \in \sR^{f_s \times h \cdot w \times d}$, which are concatenated with the video latent tokens $\vz_v^{(t)} \in \sR^{f' \times h \cdot w \times d}$ along the frame dimension:
\[
    \mathcal{T}_{total} = [\vz_v^{(t)}; \vz_s] \in \sR^{(f' + f_s) \times h \cdot w \times d}. \tag{2}
\]

This design enables anchor views to guide world synthesis in-context, without requiring architectural modifications to the base model. We further employ 3D RoPE~\cite{su2024roformer} to differentiate anchor views by assigning them distinct frame-axis positions in the positional embedding space.

\input{tex/figs/train}


\paragraph{View Pose Injection.}
Since each view corresponds to a specific 3D location in the world, its spatial pose is essential for grounding the customized content. 
We therefore inject pose information for both generated video frames and anchor views. 
The camera poses are encoded into embeddings $\vz_{pose} \in \sR^{(f' + f_s) \times 1 \times d}$ and spatially broadcast to match the latent resolution, yielding $\vz_{pose} \in \sR^{(f' + f_s) \times h \cdot w \times d}$. 
Before the self-attention layers, the pose embeddings are added to the visual tokens:
\[
    \mathcal{T}_{total} = \mathcal{T}_{total} + \vz_{pose}. \tag{3}
\]

By coupling visual tokens with spatial poses, the model can distinguish anchor views located at different positions and associate the generated egocentric trajectory with the correct local constraints. 
% This pose-aware formulation helps maintain scene consistency when the user approaches, revisits, or moves across anchored regions.

\paragraph{Text-Driven Anchor-View Evolution.}
To enable dynamic world customization, each anchor view is paired with a localized evolution description $\vt_i$ that specifies its temporal scene changes.
We inject these descriptions through cross-attention, leveraging the semantic priors of the pre-trained video model.
To preserve the locality of dynamic instructions, we restrict the interaction between text prompts and visual tokens using an attention mask. 
For a text prompt $\vt_j$, its text keys are visible only to the generated video tokens and the corresponding anchor-view tokens $\vz_s^{(j)}$:
\[
    \mathcal{M}(q, k_j) = 
    \begin{cases} 
    0, & \text{if } q \in \vz_v \text{ or } q \in \vz_s^{(j)}, \\
    -\infty, & \text{otherwise}.
    \end{cases} \tag{4}
\]
This masked cross-attention enables anchor-specific text control, allowing local scene states to evolve over time while reducing interference across different anchor views.

\subsection{Progressive Multi-Stage Training Strategy}\label{sec:3.3}

To progressively equip the model with egocentric human action control and evolvable anchor-view customization, we adopt a multi-stage training strategy, as illustrated in Figure~\ref{fig:train}. 
%
\textbf{Stage I \& II: Hybrid-View Action Control Training.}
We train the model to learn action-conditioned generation from hybrid viewpoints, where TPV videos provide complete full-body motion supervision.
% and FPV videos adapt the model to egocentric simulation.
In Stage I, the model is trained on large-scale third-person videos, where the camera parameters represent external observation viewpoints.
In Stage II, we then adapt the model to first-person videos by aligning the camera trajectory with the head pose of the character.
%
\textbf{Stage III \& IV: Evolvable Anchor-View Customization Training.}
After establishing action controllability, we train the model to incorporate anchor-view priors for world customization.
In Stage III, we train the model on static scenes to learn pose-aware anchor-view conditioning for consistent egocentric roaming.
In Stage IV, we mix in dynamic data with evolution descriptions to model text-driven local state changes.






\subsection{Hybrid-View Human Action Control}\label{sec:3.1}

\paragraph{Enhanced Egocentric Action Control via hybrid views.}
Human action contains rich spatial and interaction cues: the root trajectory determines global navigation, the limbs indicate potential interactions with the surrounding scene, and the head motion induces the egocentric viewpoint.
However, in first-person videos, most body parts are often outside the camera field of view, making direct supervision of full-body action control sparse and incomplete.
To overcome this limitation, we introduce third-person view (TPV) videos as auxiliary training data, where the full human body and its interactions with the surrounding scene are explicitly visible.
These videos provide rich interaction context and complete motion supervision, helping the model learn stronger spatial grounding between human motion and scene responses.
To support joint training on both TPV and first-person view (FPV) data within a unified framework, we formulate action conditioning in a projection-based manner.
Specifically, we represent the action condition by combining the full-body motion sequence with the camera trajectory, allowing the model to project 3D human motion into 2D visual observations under arbitrary viewpoints.
We first pre-train the model on large-scale and diverse TPV videos, where the camera parameters correspond to the external observation viewpoint, enabling the model to acquire projection knowledge and human-scene interaction priors.
Then, we adapt the model to egocentric simulation by aligning the camera parameters with the human head perspective in FPV data.
This design enables more accurate human-action control and stronger spatial pose awareness.

\paragraph{Spatial Pose Attention.} 
Inspired by prior work~\cite{fu20243dtrajmaster, li2025adaviewplanner}, we inject the pose conditions through a spatial pose attention mechanism. 
Specifically, a motion encoder first projects the input motion sequence $M \in \sR^{f \times k \times 6}$ into a latent embedding $\vz_m \in \sR^{f' \times k \times d}$, where $d$ is the model's hidden dimension. To ensure temporal alignment with the VAE-encoded~\cite{kingma2013auto} video latents, we employ temporal downsampling to match the temporal resolution $f'$. Analogously, a camera encoder processes the camera pose sequence $C \in \sR^{f \times 3 \times 4}$ into $\vz_c \in \sR^{f' \times 1 \times d}$, where the camera pose can represent either a third-person observation viewpoint or the first-person head viewpoint.

To exploit the inherent frame-wise correspondence between motion and video tokens, we concatenate the video tokens $\vz_v^{(t)}$ with the human motion tokens $\vz_m$ and camera pose tokens $\vz_c$ along the spatial dimension. This unified sequence is then processed by the spatial self-attention block:

\[
\begin{gathered}
\mT = [\vz_v^{(t)}; \vz_m; \vz_c] \in \sR^{f' \times (h \cdot w + k + 1) \times d}, \\
\vz_v^{(t)} = \vz_v^{(t)} + \text{Truncate}\left(\text{Attn}(\mW_Q \cdot \mT, \mW_K \cdot \mT, \mW_V \cdot \mT)\right)
\end{gathered}
\tag{1}
\]

The Truncate operator discards the auxiliary pose tokens, retaining only the updated video features.



\subsection{Evolvable Anchor-View Customization}\label{sec:3.2}

To enable evolvable world customization, we represent the environment with a set of anchor views. 
Each anchor view provides three types of localized world priors: an RGB image for visual appearance, a 3D pose for spatial grounding, and an evolution prompt for temporal state evolution. 

\paragraph{In-Context Anchor-View Priors.}
To incorporate anchor-view image priors while preserving the generative capability of the pre-trained video model, we adopt an in-context conditioning strategy~\cite{huang2026cinescene, ju2025fulldit, ye2025unic}. 
Specifically, the images of anchor views are encoded into latent tokens $\vz_s \in \sR^{f_s \times h \cdot w \times d}$, which are concatenated with the video latent tokens $\vz_v^{(t)} \in \sR^{f' \times h \cdot w \times d}$ along the frame dimension:
\[
    \mathcal{T}_{total} = [\vz_v^{(t)}; \vz_s] \in \sR^{(f' + f_s) \times h \cdot w \times d}. \tag{2}
\]

This design enables anchor views to guide world synthesis in-context, without requiring architectural modifications to the base model. We further employ 3D RoPE~\cite{su2024roformer} to differentiate anchor views by assigning them distinct frame-axis positions in the positional embedding space.

\input{tex/figs/train}


\paragraph{View Pose Injection.}
Since each view corresponds to a specific 3D location in the world, its spatial pose is essential for grounding the customized content. 
We therefore inject pose information for both generated video frames and anchor views. 
The camera poses are encoded into embeddings $\vz_{pose} \in \sR^{(f' + f_s) \times 1 \times d}$ and spatially broadcast to match the latent resolution, yielding $\vz_{pose} \in \sR^{(f' + f_s) \times h \cdot w \times d}$. 
Before the self-attention layers, the pose embeddings are added to the visual tokens:
\[
    \mathcal{T}_{total} = \mathcal{T}_{total} + \vz_{pose}. \tag{3}
\]

By coupling visual tokens with spatial poses, the model can distinguish anchor views located at different positions and associate the generated egocentric trajectory with the correct local constraints. 
% This pose-aware formulation helps maintain scene consistency when the user approaches, revisits, or moves across anchored regions.

\paragraph{Text-Driven Anchor-View Evolution.}
To enable dynamic world customization, each anchor view is paired with a localized evolution description $\vt_i$ that specifies its temporal scene changes.
We inject these descriptions through cross-attention, leveraging the semantic priors of the pre-trained video model.
To preserve the locality of dynamic instructions, we restrict the interaction between text prompts and visual tokens using an attention mask. 
For a text prompt $\vt_j$, its text keys are visible only to the generated video tokens and the corresponding anchor-view tokens $\vz_s^{(j)}$:
\[
    \mathcal{M}(q, k_j) = 
    \begin{cases} 
    0, & \text{if } q \in \vz_v \text{ or } q \in \vz_s^{(j)}, \\
    -\infty, & \text{otherwise}.
    \end{cases} \tag{4}
\]
This masked cross-attention enables anchor-specific text control, allowing local scene states to evolve over time while reducing interference across different anchor views.

\subsection{Progressive Multi-Stage Training Strategy}\label{sec:3.3}

To progressively equip the model with egocentric human action control and evolvable anchor-view customization, we adopt a multi-stage training strategy, as illustrated in Figure~\ref{fig:train}. 
%
\textbf{Stage I \& II: Hybrid-View Action Control Training.}
We train the model to learn action-conditioned generation from hybrid viewpoints, where TPV videos provide complete full-body motion supervision.
% and FPV videos adapt the model to egocentric simulation.
In Stage I, the model is trained on large-scale third-person videos, where the camera parameters represent external observation viewpoints.
In Stage II, we then adapt the model to first-person videos by aligning the camera trajectory with the head pose of the character.
%
\textbf{Stage III \& IV: Evolvable Anchor-View Customization Training.}
After establishing action controllability, we train the model to incorporate anchor-view priors for world customization.
In Stage III, we train the model on static scenes to learn pose-aware anchor-view conditioning for consistent egocentric roaming.
In Stage IV, we mix in dynamic data with evolution descriptions to model text-driven local state changes.






\subsection{Hybrid-View Human Action Control}\label{sec:3.1}

\paragraph{Enhanced Egocentric Action Control via hybrid views.}
Human action contains rich spatial and interaction cues: the root trajectory determines global navigation, the limbs indicate potential interactions with the surrounding scene, and the head motion induces the egocentric viewpoint.
However, in first-person videos, most body parts are often outside the camera field of view, making direct supervision of full-body action control sparse and incomplete.
To overcome this limitation, we introduce third-person view (TPV) videos as auxiliary training data, where the full human body and its interactions with the surrounding scene are explicitly visible.
These videos provide rich interaction context and complete motion supervision, helping the model learn stronger spatial grounding between human motion and scene responses.
To support joint training on both TPV and first-person view (FPV) data within a unified framework, we formulate action conditioning in a projection-based manner.
Specifically, we represent the action condition by combining the full-body motion sequence with the camera trajectory, allowing the model to project 3D human motion into 2D visual observations under arbitrary viewpoints.
We first pre-train the model on large-scale and diverse TPV videos, where the camera parameters correspond to the external observation viewpoint, enabling the model to acquire projection knowledge and human-scene interaction priors.
Then, we adapt the model to egocentric simulation by aligning the camera parameters with the human head perspective in FPV data.
This design enables more accurate human-action control and stronger spatial pose awareness.

\paragraph{Spatial Pose Attention.} 
Inspired by prior work~\cite{fu20243dtrajmaster, li2025adaviewplanner}, we inject the pose conditions through a spatial pose attention mechanism. 
Specifically, a motion encoder first projects the input motion sequence $M \in \sR^{f \times k \times 6}$ into a latent embedding $\vz_m \in \sR^{f' \times k \times d}$, where $d$ is the model's hidden dimension. To ensure temporal alignment with the VAE-encoded~\cite{kingma2013auto} video latents, we employ temporal downsampling to match the temporal resolution $f'$. Analogously, a camera encoder processes the camera pose sequence $C \in \sR^{f \times 3 \times 4}$ into $\vz_c \in \sR^{f' \times 1 \times d}$, where the camera pose can represent either a third-person observation viewpoint or the first-person head viewpoint.

To exploit the inherent frame-wise correspondence between motion and video tokens, we concatenate the video tokens $\vz_v^{(t)}$ with the human motion tokens $\vz_m$ and camera pose tokens $\vz_c$ along the spatial dimension. This unified sequence is then processed by the spatial self-attention block:

\[
\begin{gathered}
\mT = [\vz_v^{(t)}; \vz_m; \vz_c] \in \sR^{f' \times (h \cdot w + k + 1) \times d}, \\
\vz_v^{(t)} = \vz_v^{(t)} + \text{Truncate}\left(\text{Attn}(\mW_Q \cdot \mT, \mW_K \cdot \mT, \mW_V \cdot \mT)\right)
\end{gathered}
\tag{1}
\]

The Truncate operator discards the auxiliary pose tokens, retaining only the updated video features.



\subsection{Evolvable Anchor-View Customization}\label{sec:3.2}

To enable evolvable world customization, we represent the environment with a set of anchor views. 
Each anchor view provides three types of localized world priors: an RGB image for visual appearance, a 3D pose for spatial grounding, and an evolution prompt for temporal state evolution. 

\paragraph{In-Context Anchor-View Priors.}
To incorporate anchor-view image priors while preserving the generative capability of the pre-trained video model, we adopt an in-context conditioning strategy~\cite{huang2026cinescene, ju2025fulldit, ye2025unic}. 
Specifically, the images of anchor views are encoded into latent tokens $\vz_s \in \sR^{f_s \times h \cdot w \times d}$, which are concatenated with the video latent tokens $\vz_v^{(t)} \in \sR^{f' \times h \cdot w \times d}$ along the frame dimension:
\[
    \mathcal{T}_{total} = [\vz_v^{(t)}; \vz_s] \in \sR^{(f' + f_s) \times h \cdot w \times d}. \tag{2}
\]

This design enables anchor views to guide world synthesis in-context, without requiring architectural modifications to the base model. We further employ 3D RoPE~\cite{su2024roformer} to differentiate anchor views by assigning them distinct frame-axis positions in the positional embedding space.

\input{tex/figs/train}


\paragraph{View Pose Injection.}
Since each view corresponds to a specific 3D location in the world, its spatial pose is essential for grounding the customized content. 
We therefore inject pose information for both generated video frames and anchor views. 
The camera poses are encoded into embeddings $\vz_{pose} \in \sR^{(f' + f_s) \times 1 \times d}$ and spatially broadcast to match the latent resolution, yielding $\vz_{pose} \in \sR^{(f' + f_s) \times h \cdot w \times d}$. 
Before the self-attention layers, the pose embeddings are added to the visual tokens:
\[
    \mathcal{T}_{total} = \mathcal{T}_{total} + \vz_{pose}. \tag{3}
\]

By coupling visual tokens with spatial poses, the model can distinguish anchor views located at different positions and associate the generated egocentric trajectory with the correct local constraints. 
% This pose-aware formulation helps maintain scene consistency when the user approaches, revisits, or moves across anchored regions.

\paragraph{Text-Driven Anchor-View Evolution.}
To enable dynamic world customization, each anchor view is paired with a localized evolution description $\vt_i$ that specifies its temporal scene changes.
We inject these descriptions through cross-attention, leveraging the semantic priors of the pre-trained video model.
To preserve the locality of dynamic instructions, we restrict the interaction between text prompts and visual tokens using an attention mask. 
For a text prompt $\vt_j$, its text keys are visible only to the generated video tokens and the corresponding anchor-view tokens $\vz_s^{(j)}$:
\[
    \mathcal{M}(q, k_j) = 
    \begin{cases} 
    0, & \text{if } q \in \vz_v \text{ or } q \in \vz_s^{(j)}, \\
    -\infty, & \text{otherwise}.
    \end{cases} \tag{4}
\]
This masked cross-attention enables anchor-specific text control, allowing local scene states to evolve over time while reducing interference across different anchor views.

\subsection{Progressive Multi-Stage Training Strategy}\label{sec:3.3}

To progressively equip the model with egocentric human action control and evolvable anchor-view customization, we adopt a multi-stage training strategy, as illustrated in Figure~\ref{fig:train}. 
%
\textbf{Stage I \& II: Hybrid-View Action Control Training.}
We train the model to learn action-conditioned generation from hybrid viewpoints, where TPV videos provide complete full-body motion supervision.
% and FPV videos adapt the model to egocentric simulation.
In Stage I, the model is trained on large-scale third-person videos, where the camera parameters represent external observation viewpoints.
In Stage II, we then adapt the model to first-person videos by aligning the camera trajectory with the head pose of the character.
%
\textbf{Stage III \& IV: Evolvable Anchor-View Customization Training.}
After establishing action controllability, we train the model to incorporate anchor-view priors for world customization.
In Stage III, we train the model on static scenes to learn pose-aware anchor-view conditioning for consistent egocentric roaming.
In Stage IV, we mix in dynamic data with evolution descriptions to model text-driven local state changes.




